\section{Square grids}\label{sec:square}

On the $n\times n$ grid the row bound is of order $n^2$ and is far from the truth. For
$n\ge3$, Lemma~\ref{lem:girth} and Neuwirth's bound for bipartite graphs of girth $8$ with
$n$ vertices on each side~\cite[Corollary~16]{neuwirth2001} give
$f(n,n)<n^{4/3}+\tfrac23\,n$. The incidence graphs of generalized quadrangles show that
this is the right order for axis-aligned rectangles alone~\cite{benson1966}. With tilted
rectangles we only have the following lower bound.

\begin{remark}\label{rem:square}
There is a constant $c>0$ such that $f(n,n)\ge c\,n^{6/5}(1+\ln n)^{-2/5}$ for every
$n\ge2$. In particular $f(n,n)-2n\to\infty$. We only sketch the proof, which is the
deletion method~\cite[Section~3.2]{alonspencer2016}. Put $\Lambda=1+\ln n$. By
Lemmas~\ref{lem:diff}, \ref{lem:threecorners}, and~\ref{lem:compress}, a blackout is valid
if it contains no set of the following 2 families. The family $\mathcal A$ consists of the
corner sets of rectangles and of the supports of mixed pairs, which have 4 points. The
family $\mathcal B$ consists of the sets $C(R_1)\triangle C(R_2)$ for rectangles with
exactly 1 common corner, which have 6 points. A count of the sides at a point shows that
at most $8n^2\Lambda$ rectangles have a corner at a given point. This gives
$|\mathcal B|\le32n^6\Lambda^2$, and with the divisor bound it gives
$|\mathcal A|\le Kn^{4+1/10}\Lambda$. Choose each point independently with probability
$a\,n^{-4/5}\Lambda^{-2/5}$, where $a=192^{-1/5}$, and delete 1 point from each member of
$\mathcal A\cup\mathcal B$ whose points are all chosen. The expected number of points that
remain is at least $(\frac56a-o(1))\,n^{6/5}\Lambda^{-2/5}$. The exponent $6/5$ comes from
the roughly $n^6$ members of $\mathcal B$.
\end{remark}
